% Architecture du cluster MyApp — Document de partage
% Compilation : pdflatex architecture-cluster-myapp.tex

\documentclass[11pt,a4paper]{article}
\usepackage[utf8]{inputenc}
\usepackage[T1]{fontenc}
\usepackage[french]{babel}
\usepackage{geometry}
\usepackage{hyperref}
\usepackage{booktabs}
\usepackage{tabularx}
\usepackage{enumitem}
\usepackage{listings}
\usepackage{xcolor}
\usepackage{graphicx}
\usepackage{fancyhdr}
\usepackage{titlesec}

\geometry{margin=2.5cm}
\hypersetup{colorlinks=true, linkcolor=blue, urlcolor=blue}
\setlength{\parindent}{0pt}
\setlength{\parskip}{0.6em}

% Style des listings
\lstset{
  basicstyle=\ttfamily\small,
  breaklines=true,
  frame=single,
  backgroundcolor=\color{gray!10},
  numbers=left,
  numberstyle=\tiny
}

\title{\textbf{Architecture du cluster MyApp} \\ \large État actuel, responsabilités DevOps \& développeurs}
\author{Équipe MyApp}
\date{\today}

\begin{document}

\maketitle
\tableofcontents
\newpage

%==============================================================================
\section{Introduction}
%==============================================================================

Ce document décrit l'architecture du cluster Kubernetes MyApp, son état actuel, et répartit les responsabilités entre l'ingénieur DevOps et les développeurs de chaque repository. Il est destiné à être partagé avec l'équipe pour assurer une vision commune.

\textbf{Stack technique :}
\begin{itemize}[noitemsep]
  \item Cluster k3s sur GCP (2 nœuds)
  \item Namespace : \texttt{myapp}
  \item Ingress : Traefik (intégré k3s)
  \item Accès : \texttt{http://203.0.113.10}
\end{itemize}

%==============================================================================
\section{Architecture du cluster}
%==============================================================================

\subsection{Schéma des nœuds}

Le cluster comprend \textbf{2 nœuds} :

\begin{center}
\begin{tabular}{@{}lllll@{}}
\toprule
\textbf{Nœud} & \textbf{Rôle} & \textbf{IP interne} & \textbf{CPU} & \textbf{RAM} \\
\midrule
backend-vm  & control-plane, master & 10.0.0.11 & 2 cores & $\sim$4 Go \\
frontend-vm & worker                 & 10.0.0.12 & 2 cores & $\sim$4 Go \\
\bottomrule
\end{tabular}
\end{center}

\subsection{Répartition des services par nœud}

Les applications sont déployées en \textbf{haute disponibilité} (2 replicas) avec répartition sur les deux nœuds via \texttt{topologySpreadConstraints}.

\begin{center}
\begin{tabular}{@{}p{5cm}p{6cm}@{}}
\toprule
\textbf{backend-vm} & \textbf{frontend-vm} \\
\midrule
PostgreSQL (1 replica) & Consul (répliqué) \\
Redis (1 replica)     & \\
Consul               & \\
RabbitMQ (1 replica) & \\
\midrule
api-gateway $\times$2     & api-gateway $\times$2 \\
user-management $\times$2 & user-management $\times$2 \\
payment-service $\times$2 & payment-service $\times$2 \\
crm-client $\times$2      & crm-client $\times$2 \\
predictions-intake $\times$2 & predictions-intake $\times$2 \\
kpi-dashboard $\times$2   & kpi-dashboard $\times$2 \\
chatbot $\times$2         & chatbot $\times$2 \\
frontend $\times$2        & frontend $\times$2 \\
\bottomrule
\end{tabular}
\end{center}

\subsection{Flux réseau et Ingress}

\begin{itemize}[noitemsep]
  \item \textbf{IP publique} : \texttt{203.0.113.10}
  \item \textbf{Ingress principal} (\texttt{ingress-ip}) : \\
  \begin{tabular}{ll}
    \texttt{/api} & $\rightarrow$ api-gateway (8888) \\
    \texttt{/}    & $\rightarrow$ frontend (3000)   \\
  \end{tabular}
  \item \textbf{CORS} : api-gateway accepte \texttt{app.localhost}, \texttt{203.0.113.10}, \texttt{localhost:3000}
  \item Les autres services (CRM, KPI, payment, etc.) sont atteignables via l’api-gateway et Consul
\end{itemize}

%==============================================================================
\section{État actuel}
%==============================================================================

\subsection{Ce qui est en place}

\begin{itemize}
  \item Cluster k3s opérationnel
  \item HA : 2 replicas par application, répartis sur les 2 nœuds
  \item RabbitMQ : 1 replica sur backend-vm, nettoyage des pods evicted
  \item CORS : configuré pour accès par IP
  \item CronJobs : nettoyage disque sur backend-vm et frontend-vm
  \item Secrets : postgres-credentials (postgres/postgres)
\end{itemize}

\subsection{Ce qui reste à configurer}

\begin{enumerate}
  \item \textbf{ghcr-secret} : à créer manuellement (GitHub PAT)
  \item \textbf{Google OAuth} : Client ID, Secret, Redirect URI
  \item \textbf{Stripe} : clés API et webhook (si paiements utilisés)
  \item \textbf{JWT\_SECRET} : secret partagé api-gateway / user-management
  \item \textbf{OAuth redirect} : \texttt{CustomOAuth2SuccessHandler} utilise \texttt{localhost:3000} en dur
  \item \textbf{Postgres} : mot de passe plus fort en production
\end{enumerate}

%==============================================================================
\section{Rôle de l'ingénieur DevOps}
%==============================================================================

\subsection{Responsabilités}

\begin{enumerate}
  \item \textbf{Infrastructure}
  \begin{itemize}[noitemsep]
    \item Maintien du cluster k3s (backend-vm, frontend-vm)
    \item Vérification des nœuds et du placement des pods
  \end{itemize}

  \item \textbf{Secrets à créer manuellement}
  \begin{lstlisting}
kubectl create secret docker-registry ghcr-secret \
  --namespace=myapp \
  --docker-server=ghcr.io \
  --docker-username=TON_GITHUB_USERNAME \
  --docker-password=TON_GITHUB_PAT
  \end{lstlisting}

  \item \textbf{Postgres} : en prod, remplacer postgres/postgres par un mot de passe fort (ou Sealed Secrets / Vault)

  \item \textbf{Configuration réseau}
  \begin{itemize}[noitemsep]
    \item Si l’IP publique change : mettre à jour \texttt{apps/api-gateway/deployment.yaml} \\
    (FRONTEND\_ORIGIN, SPRING\_APPLICATION\_JSON)
    \item Vérifier les Ingress et le routage
  \end{itemize}

  \item \textbf{Déploiement}
  \begin{lstlisting}
kubectl apply -k deploy/k8s/
  \end{lstlisting}

  \item \textbf{Monitoring}
  \begin{itemize}[noitemsep]
    \item Vérifier les pods : \texttt{kubectl get pods -n myapp -o wide}
    \item Vérifier les Ingress : \texttt{kubectl get ingress -n myapp}
    \item Nettoyer les pods evicted si nécessaire
  \end{itemize}
\end{enumerate}

%==============================================================================
\section{Rôle de chaque développeur (par repository)}
%==============================================================================

\subsection{Backend User-Management (\texttt{Backend-User-management-and-subscripption})}

\begin{itemize}
  \item \textbf{JWT\_SECRET} : fournir une valeur (32+ bytes) à injecter en k8s ; doit être la même que pour api-gateway
  \item \textbf{GOOGLE\_CLIENT\_ID} et \textbf{GOOGLE\_CLIENT\_SECRET} : fournis par le dev après création de l’app Google Cloud
  \item \textbf{GOOGLE\_REDIRECT\_URI} : \texttt{http://203.0.113.10/api/auth/oauth2/code/google} (ou URL prod)
  \item \textbf{Code à modifier} : \texttt{CustomOAuth2SuccessHandler.java} ligne $\sim$155 : remplacer \texttt{http://localhost:3000} par une variable d’environnement \texttt{FRONTEND\_URL}
\end{itemize}

\subsection{Backend Payment-Service (\texttt{Backend-payment-service})}

\begin{itemize}
  \item \textbf{STRIPE\_API\_KEY} et \textbf{STRIPE\_WEBHOOK\_SECRET} : fournis par le dev Stripe
  \item \textbf{stripe.success.url} et \textbf{stripe.cancel.url} : configurer les URLs frontend (ex. \texttt{http://203.0.113.10/dashboard})
  \item \textbf{RABBITMQ\_USERNAME} / \textbf{RABBITMQ\_PASSWORD} : ajouter dans le deployment k8s (guest/guest pour le cluster actuel)
  \item S’assurer que le service utilise bien \texttt{spring.rabbitmq.host/port} pour le cluster k8s (pas CloudAMQP uniquement)
\end{itemize}

\subsection{Backend API-Gateway (\texttt{Backend-api-gateway})}

\begin{itemize}
  \item \textbf{JWT\_SECRET} : injecter la même valeur que user-management
  \item CORS : déjà configuré via \texttt{SPRING\_APPLICATION\_JSON} dans le deployment
\end{itemize}

\subsection{Frontend (\texttt{Front-end})}

\begin{itemize}
  \item \textbf{NEXT\_PUBLIC\_API\_URL} : \texttt{/api} (déjà correct pour accès relatif)
  \item Vérifier que l’OAuth callback (\texttt{/auth/oauth2/code/google}) correspond à l’URL configurée côté user-management
\end{itemize}

\subsection{Autres services (KPI-Dashboard, CRM-Client, Predictions-Intake, Chatbot)}

\begin{itemize}
  \item S’enregistrer dans Consul avec le nom attendu par l’api-gateway :
  \begin{itemize}[noitemsep]
    \item user-management $\rightarrow$ \texttt{user-service}
    \item crm-client $\rightarrow$ \texttt{crmservice}
    \item payment-service $\rightarrow$ \texttt{payment-service}
    \item predictions-intake $\rightarrow$ \texttt{prediction-intake-service}
    \item kpi-dashboard $\rightarrow$ \texttt{kpi-service}
  \end{itemize}
  \item Fournir des images Docker publiées sur ghcr.io
  \item Documenter les variables d’environnement requises
\end{itemize}

%==============================================================================
\section{Checklist globale}
%==============================================================================

\subsection{Priorité critique}

\begin{enumerate}
  \item[$\square$] Créer \texttt{ghcr-secret} (DevOps)
  \item[$\square$] Configurer Google OAuth (Client ID, Secret, Redirect URI) — dev user-management
  \item[$\square$] Rendre \texttt{FRONTEND\_URL} configurable dans \texttt{CustomOAuth2SuccessHandler} — dev user-management
  \item[$\square$] Configurer Stripe (clés + URLs) si paiements utilisés — dev payment-service
\end{enumerate}

\subsection{Priorité importante (sécurité)}

\begin{enumerate}
  \item[$\square$] JWT\_SECRET partagé (api-gateway + user-management)
  \item[$\square$] Mot de passe Postgres renforcé en prod (DevOps)
\end{enumerate}

\subsection{Priorité optionnelle}

\begin{enumerate}
  \item[$\square$] HTTPS (certificats + config Traefik)
  \item[$\square$] Domaine personnalisé au lieu de l’IP
\end{enumerate}

%==============================================================================
\section{Références}
%==============================================================================

\begin{itemize}
  \item \texttt{deploy/k8s/README.md} — structure des manifests
  \item \texttt{deploy/k8s/DEPLOYMENT.md} — ordre de déploiement
  \item \texttt{deploy/k8s/CHECKLIST.md} — checklist détaillée
  \item Vérifications rapides :
  \begin{lstlisting}
kubectl get pods -n myapp
kubectl get ingress -n myapp
curl -s http://203.0.113.10/
curl -s http://203.0.113.10/api/actuator/health
  \end{lstlisting}
\end{itemize}

\vspace{1em}
\hrule
\vspace{0.5em}
\textit{Document généré pour partage avec l’équipe MyApp.}

\end{document}
